\section{Related Work}
\label{sec:related}

\subsection{The SESGO benchmark}
\label{sec:related-sesgo}

\textsc{SESGO}~\citep{robles2025sesgo} is a Spanish, Latin-American-grounded social-bias
benchmark in the BBQ ambiguous/disambiguated tradition, spanning racism,
xenophobia, classism, and gender. Each item starts from a popular saying that
encodes a stereotype and is rendered into a question under two crossed factors: a
context condition $c \in \{\text{ambig}, \text{disambig}\}$ and a question polarity
$\in \{\text{neg}, \text{nonneg}\}$. The answer options realize three roles: the
stereotyped \emph{Target} group ($t$), an \emph{Other} group ($o$), and
\emph{Unknown} ($u$). On an \emph{ambiguous} item the text names no specific
person, so the gold answer is $u$ (abstain); on a \emph{disambiguated} item the text
specifies a group, and the gold is that group. We write $F(\text{Target})$ and
$F(\text{Other})$ for the share of \emph{incorrect} answers that harm the
marginalized Target and the Other group respectively, and define the signed
\emph{bias alignment} as $F(\text{Target}) - F(\text{Other})$, positive when errors
lean toward the stereotyped group. \textsc{SESGO} summarizes a model's behavior on a
condition with a single
\begin{equation}
  \mathrm{bias\_score}
    \;=\; \sigma \cdot \sqrt{\,(1 - \mathrm{acc})^2 + \big(F(\text{Target}) - F(\text{Other})\big)^2\,},
  \label{eq:bias-score}
\end{equation}
which is large when the model is both inaccurate and one-sidedly biased ($\sigma$ a
fixed sign/scale factor; \citealp{robles2025sesgo}). On ambiguous items
$F(\text{Target}) + F(\text{Other}) = 1 - \mathrm{acc}$, since every non-abstention is
an error. Our extension is along \emph{model scale}: we replicate this scoring across
open-weight models of varying size.


\subsection{Chain-of-thought dynamics}
\label{sec:related-cot}

Reasoning models commit to an answer somewhere along their chain of thought, and
the forking-paths method~\citep{forkingpaths,forkingpaths2} localizes \emph{where}.
Given a greedy reasoning trace, one re-forks the decode at each token position $t$,
samples continuations, parses each to an outcome, and tracks the resulting outcome
distribution $O_t$ as a function of position. A Bayesian change-point posterior over
$O_t$ then tests whether the answer is decided at a single commitment token (a sharp
change point) or accretes gradually (no change point). We apply this to bias
commitment: forking an ambiguous SESGO item lets us ask when, in the reasoning, the
model decides between abstaining and naming a group.


%%%%%%%%%%%%%%%%%%%%%%%%%%%%%%%%%%%%%%%%%%%%%%%%%%%%%%%%%%%%%%%%%%%%%%%%%%
